% !TeX root = ../../Book2.tex
% 纲要标题：### 15.4 边界的领状邻域
% 纲要位置：几何学计划纲要.md，第 1597 行。
% 正文状态：初稿；尚未正式审核。

% 纲要内容开始
% 
% * 向内横截向量场与局部流
% * 边界邻域的整体拼接
% * 双化和相对粘合中的应用
% 
% 纲要内容结束

\section{边界的领状邻域}
\label{g2:sec:1504}

在半空间图中，边界附近看起来像“边界坐标加一个非负坐标”。各张图的非负坐标并不相同，因而还不能把所有边界点统一向内推一段距离。本节用一个向内向量场把这些局部图连接起来：每条轨道从边界出发，时间成为共同的横向坐标。由此得到的领状邻域可用于双化、粘合和保持边界数据的构造。

\subsection{向内横截向量场与局部流}
\label{g2:subsec:1504:01}

设 \(M\) 为有限维、Hausdorff、第二可数的光滑带边界流形。在边界坐标 \((y,r)\)、\(r\ge0\) 中，向量 \(v\in\mathrm{T}_xM\) 严格指向内部的条件是 \(\mathrm{d}r_x(v)>0\)。第一卷第11章的边界坐标变换说明，另一边界定义坐标沿边界的微分是 \(\mathrm{d}r\) 的正倍数，所以这一条件不依赖所选图。它也说明同一点的向内向量组成凸集，正的加权平均仍向内。

\begin{lemma}\label{g2:lem:boundary-defining-function}
	每个光滑带边界流形 \(M\) 都有光滑函数 \(r:M\to[0,\infty)\)，满足 \(r^{-1}(0)=\partial M\) 且 \(\mathrm{d}r\ne0\) 于 \(\partial M\)。在边界的某个开邻域上，可选择光滑向量场 \(X\)，使 \(Xr=1\)。
\end{lemma}
\begin{proof}
	取由边界坐标图和内点坐标图组成的局部有限覆盖，并取从属光滑单位分解 \(\chi_i\)。边界图上记最后一个坐标为 \(r_i\)，内点图上令 \(r_i=1\)，定义 \(r=\sum_i\chi_i r_i\)，各项在支集外补零。每个内点的所有非零项都为正，故 \(r>0\)；在边界点，来自内点图的项在该点附近消失，其余 \(r_i\) 为零，所以 \(r=0\)。边界上
	\[
		\mathrm{d}r=\sum_i\chi_i\,\mathrm{d}r_i.
	\]
	各个 \(\mathrm{d}r_i\) 是同一个向内法向余向量的正倍数，且权重之和为1，故和非零。

	在各边界图中取 \(\partial_{r_i}\)，用单位分解加权拼成边界附近的光滑向量场 \(Y\)。其边界值严格向内，因此 \(Yr>0\) 在边界的某开邻域 \(U\) 上成立。令 \(X=Y/(Yr)\) 于 \(U\)，便有 \(Xr=1\)。这里使用的是第一卷第12章第3节已经证明的局部有限单位分解；局部有限性保证求和及逐项求导有意义。
\end{proof}

\(r\) 称为一个\term[bianjiedingyihanshu]{边界定义函数}[boundary defining function]。条件 \(\mathrm{d}r\ne0\) 保证它真正测量横向变化；仅有同一零点集的函数 \(r^2\) 不具有这个性质。归一化 \(Xr=1\) 的作用也很具体：沿从边界出发的积分曲线，\(r\) 的值恰好就是已经经过的时间。

\ProseParagraphBreak

在边界处不能直接要求 \(X\) 生成保持 \(M\) 的双侧流。向内的速度在负时间会走到半空间外。我们在边界坐标中把光滑分量局部延拓到 Euclid 开集，应用第一卷第13章第1节的初值存在、唯一和光滑依赖定理，再保留非负时间。因 \(r(0)=0\)、\(r'=1\)，所得小段正时间曲线位于内部。不同局部延拓在半空间内给出同一个方程，解的唯一性使这些正时间解一致。记其为
\[
	F(x,t)=\Phi_t(x),\qquad x\in\partial M,\quad 0\le t<\varepsilon_x.
\]
在 \(t=0\) 处，\(\mathrm{d}F\) 把边界切向量保持为原向量，并把 \(\partial_t\) 送到 \(X_x\)。这个向量不切于边界，因此上述导数为同构，局部逆函数论证给出边界附近的半空间坐标。

\subsection{边界邻域的整体拼接}
\label{g2:subsec:1504:02}

\begin{definition}\label{g2:def:collar}
	设 \(M\) 为光滑带边界流形。若光滑嵌入 \(c:\partial M\times[0,1)\to M\) 满足 \(c(x,0)=x\)，且像是 \(\partial M\) 在 \(M\) 中的开邻域，就称它为一个\term[lingzhuanglinyu]{领状邻域}[collar neighborhood]的参数化，简称领圈。
\end{definition}

\begin{theorem}[领状邻域定理]\label{g2:thm:collar-neighborhood}
	每个有限维、Hausdorff、第二可数的光滑带边界流形都存在领圈。边界可以非紧；领圈的像可取在任何预先给定的边界开邻域内。
\end{theorem}
\begin{proof}
	在前面构造的开集 \(U\) 上使用 \(Xr=1\) 的向量场，再把 \(U\) 与指定的边界邻域相交。由局部存在与光滑依赖，每个边界点有一个边界邻域 \(V_i\) 及正数 \(a_i\)，使从 \(V_i\) 出发的解在 \(0\le t<a_i\) 上存在并留在 \(U\)。可取 \(\{V_i\}\) 局部有限。设 \(\{\psi_i\}\) 为从属的单位分解，令
	\[
		\varepsilon(x)=\frac12\sum_i\psi_i(x)a_i>0.
	\]
	对每个 \(x\)，在 \(\psi_i(x)>0\) 的有限多个指标中选 \(a_i\) 最大者，则 \(\varepsilon(x)<a_i\) 且 \(x\in V_i\)。故 \(F\) 在
	\[
		D=\{(x,t):x\in\partial M,\ 0\le t<\varepsilon(x)\}
	\]
	上有定义，光滑且满足 \(r(F(x,t))=t\)。

	若 \(F(x,t)=F(y,s)\)，对两边取 \(r\) 得 \(t=s\)。两条轨道从相同时间、相同位置反向求解，在返回边界前都位于 \(M\)，唯一性给出 \(x=y\)。所以 \(F\) 在整个 \(D\) 上单射，这一步同时排除了来自遥远边界点的相撞。

	在正时间，流的微分把 \(\mathrm{T}_x\partial M\) 与 \(X_x\) 组成的基送到终点，并保持生成方向 \(X\)，故 \(\mathrm{d}F\) 仍可逆。也可注意 \(r\circ F=t\)，再在每张 \(t\) 切片上用流的局部逆。零时间处已由半空间局部逆得到同样结论。于是 \(F\) 是单射局部微分同胚，其局部逆在交叠处一致，拼成到开像的光滑整体逆。

	最后令 \(c(x,u)=F(x,u\varepsilon(x))\)，\(0\le u<1\)。\((x,u)\mapsto(x,u\varepsilon(x))\) 是到 \(D\) 的微分同胚，因而 \(c\) 是所需领圈。
\end{proof}

非紧边界上，轨道允许走的时间随点变化；最后的重参数化把这个变化吸收到 \(\varepsilon(x)\) 中。统一参数区间 \([0,1)\) 因而不表示每个点都能沿最初向量场运动同样长的时间。相关流与领圈构造可参看\cite[第217--223页，定理9.20、9.24、9.25]{Lee2012Smooth}。

\begin{example}\label{g2:exm:disk-collar}
	闭单位球 \(D^n\) 的一个领圈为 \(c(x,t)=(1-t/2)x\)，其中 \(x\in S^{n-1}\)、\(0\le t<1\)。其像是 \(\{z:1/2<|z|\le1\}\)，在 \(D^n\) 中开。边界点被逐渐推入内部，角方向保持不变；半径 \(1/2\) 没有包含在像中，对应参数端点 \(t=1\) 未取到。
\end{example}

\subsection{双化和相对粘合中的应用}
\label{g2:subsec:1504:03}

取 \(M\) 的两份副本 \(M_+,M_-\)，将对应的边界点识别，得到\term[shuanghua]{双化}[double]
\[
	\operatorname{Dbl}(M)=(M_+\amalg M_-)/\{x_+\sim x_-:x\in\partial M\}.
\]
给定领圈 \(c\) 后，用一个有符号的参数把接缝附近写成
\[
	\Theta:\partial M\times(-1,1)\longrightarrow\operatorname{Dbl}(M),\qquad
	\Theta(x,s)=\begin{cases}
		[c(x,s)]_+,&s\ge0,\\
		[c(x,-s)]_-,&s\le0.
	\end{cases}
\]
两条表达在零处一致。它是到接缝开邻域的同胚：在两半上分别有领圈的逆，它们的第一分量在边界一致，第二分量分别取正号、负号，因而由粘合定义得到连续逆。

\begin{proposition}\label{g2:prop:smooth-boundary-gluing}
	设 \(M_1,M_2\) 为同维光滑带边界流形，\(\varphi:\partial M_1\to\partial M_2\) 为微分同胚。选定两侧领圈后，商空间 \(M_1\cup_\varphi M_2\) 具有光滑无边界流形结构，使两份 \(M_i\) 都以闭的带边界子流形嵌入其中。特别地，\(\operatorname{Dbl}(M)\) 是无边界光滑流形。只沿边界分支的并作同样粘合时，其余边界保留。
\end{proposition}
\begin{proof}
	以 \(\partial M_1\) 为共同边界，在负半边用 \(c_2(\varphi(x),-s)\) 代替上面的表达。两侧内点保留原图；接缝附近取 \(\partial M_1\) 的坐标与 \(s\) 构成乘积图。两张接缝图的变换为 \((y,s)\mapsto(h(y),s)\)，所以光滑。接缝图与内点图的交落在 \(s>0\) 或 \(s<0\)，过渡由领圈、原流形坐标及 \(\varphi\) 的光滑性保证。

	商空间还须满足分离性与可数性。每个接缝点有上述乘积邻域，两侧内点也有原来的开邻域。两个不同接缝点用边界上的不交开集及较小领圈分离；内点与接缝点用避开该内点的小领圈分离；内点之间使用原流形的 Hausdorff 性或两侧互不相交的内部。这证明 Hausdorff 性。内部图及接缝的乘积图均可取可数族，故整体第二可数。各侧在这些坐标中为闭半空间，其补集是另一侧内部，所以是所述闭嵌入。
\end{proof}

粘合光滑对象时，领圈还告诉我们应匹配哪些数据。若两边的函数在领圈内都等于同一个边界函数 \(a(x)\)，它们就粘成光滑函数，因为接缝表达为与 \(s\) 无关的 \(a(x)\)。只要求边界值相等则不够：两侧都取领圈坐标 \(u\)，接起来变成 \(|s|\)，在零处不可微。对一般函数，需要各阶法向导数按 \(u=-s\) 的符号相容；先在领圈内调整为乘积形式，是一种直接保证相容的做法。\cite[第224--225页，定理9.29]{Lee2012Smooth}给出了相同的流形粘合构造。
